% 原书第 541 页（PDF 第 554 页）
\BookPage{541}{554}
\BookHeading{盛清社会与扬州文化}{盛清社会与扬州文化}{阎崇年*}
盛清的康熙、雍正、乾隆三朝，扬州文化，位置凸显，尤为繁盛，出现高峰。本文以扬州文化中的三个史例一康熙朝编修《全唐诗集》、雍正朝蒋衡书写「十三经」和乾隆朝纂修《四库全书》为主，兼及其他，就盛清社会与扬州文化的关系，勾稽史料，略加探讨。

康熙朝编修《全唐诗》，功在当朝，影响至今。《全唐诗》在扬州刻印，既是扬州文化的一次展演，也是扬州文化的一大贡献。

《全唐诗》为康熙帝敕编。康熙帝喜欢诗，着有《圣祖仁皇帝御制诗集》初集、二集、三集、四集，共176卷，现能看到康熙帝的诗词有1147首。1写诗的人，自然爱诗，也喜读诗。康熙帝御定《全金诗》74卷；《御定四朝诗》即宋、金、元、明的诗312卷，收作者五千八百人；《御定题画诗》120卷；《御定历代赋汇》140卷、「外集」20卷、「逸句」2卷、「补遗」22卷，合计184卷；《御选唐诗》32卷等，共计722卷。

唐诗是中国诗歌史上的一个高峰，康熙帝喜欢并学习唐诗。他在文化史上做的一件盛事，就是敕编《全唐诗集》，尔后定名为《御定全唐诗》，凡九百卷。唐人诗集，从北宋以来，虽有选录的总集，却无唐诗的全集。明朝浙江海盐人胡亨震编的《唐音统签》，其汇编之唐诗，虽初具规模，却尚多遗漏。清初扬州泰兴人季振

\vfill
\begin{BookNotes}
\BookNote{}{作者任职于北京社会科学院。}
\BookNote{1.}{王志民、王泽远校注，《康熙诗词集注》，内蒙古人民出版社，呼和浩特，1993年。}
\end{BookNotes}

% 原书第 542 页（PDF 第 555 页）
\BookPage{542}{555}
\noindent 宜编辑《全唐诗》717卷，\BookRef{2}稿方成，而身故，存写本，未付梓。康熙帝亲自敕编，以内府所藏唐人诗集为基础，借鉴《唐音统签》、季振宜写本《全唐诗》，参考残碑断碣，诗笺篇什，稗史杂书，汇集唐诗四万八千九百余首，作者二千二百余人，汇成《御定全唐诗》。这部诗集，「根据诸本，一一校注，尤为周密，得此一编，而唐诗之源流、正变、始末厘然。自有总集以来，更无如是之既博且精者矣」\BookRef{3}。《全唐诗》至今仍是许多学人、特别是诗人插架的必备之书。

《全唐诗》的刊印，由时任江宁织造、通政使司通政使唐寅负责，在扬州天宁寺刻印。现中国第一历史档案馆珍藏的〈奏为《全唐诗集》定期于天宁寺开局刊刻等事〉（康熙四十四年五月初一日），记载了这件事情。清宫珍藏原档记载：

\begin{BookQuote}
江宁织造、通政使司通政使臣曹寅谨奏：臣寅恭蒙谕旨，刊刻《全唐诗集》。命词臣彭定求等九员校刊。臣寅已行文，期于五月初一日天宁寺开局，至今尚未到扬，俟其到齐校刊，谨当奏闻。又，闰四月二十三日，有翰林院庶吉士臣俞梅，赴臣寅衙门，口传上谕，命臣俞梅就近校刊《全唐诗集》。钦此。奏请圣旨。钦遵。咨行江苏巡抚臣宋荦，移咨吏部、翰林院衙门，俟刊刻完日，该衙门一并具本奏闻。康熙四十四年五月初一日。
\end{BookQuote}

\begin{BookQuote}
（朱批）知道了。\BookRef{4}
\end{BookQuote}

这一年恰为康熙帝第五次南巡，往返途经并驻跸扬州。在扬州诗局开设后，康熙帝亲自钦派江浙翰林彭定求、俞梅、汪士𬭎、汪绎、徐树本、沈三曾、杨中讷、查嗣瑮、车鼎晋、潘从律十位等，\BookRef{5}被称为「十翰林」或「十学士」。康熙帝还就书的〈凡例〉、刻板、版式、纸张、装潢、进度等做出谕旨。《全唐诗》：「于康熙四十六年（1707）四月十六日书成。谨装潢成帙，进呈圣览。（朱批）知道了。」

\vfill
\begin{BookNotes}
\BookNote{2.}{季振宜，〈全唐诗序〉，《全唐诗季振宜写本》，第1册第1页，《故宫珍本丛刊》，第621册，海南出版社，海口，2000年。}
\BookNote{3.}{《四库全书总目》，第190卷，中华书局，北京，1965年。}
\BookNote{4.}{中国第一历史档案馆、扬州市档案局（馆）编：《清宫扬州御档》，第1册，第75页，广陵书社，扬州，2009年。}
\BookNote{5.}{彭定求，《南畇老人自订年谱》，《北京图书馆藏珍本年谱丛刊》，第86册，北京图书馆出版社，北京，1999年。}
\end{BookNotes}

% 原书第 543 页（PDF 第 556 页）
\BookPage{543}{556}
\noindent 6曹寅奏称:「此皆皇上圣心独运，定为必传之书。」7在扬州编纂刻印的《全唐诗》，写刻一丝不苟，字体前后一致，「散帙编摩，订疑晰误」，⁸校勘精细，纸墨精良，装潢美观，书林典范，乃「成一代之书」，至今仍被誉为清代雕版史上的杰作。

在扬州编纂刻印《全唐诗》资费的来源，学者多认为出自「盐羡」；但有学者认为：「应肯定扬州诗局的经费来自国库。」10

在扬州除刊刻《全唐诗集》外，还刻印《佩文韵府》。先是，颜真卿的《韵海镜源》已失传，凌氏、阴氏的韵书，不够完备；至是，康熙帝敕编较前人更为完备的《佩文韵府》。在敕编《佩文韵府》全书过程中，康熙帝不仅原则指示，而且参与编审一文臣「每缮初稿，先呈御览」，阅后「点摘阙疑，…亲加批乙，宣付诸臣，再三稽考」。11随后，又有《御定韵府拾遗》120卷。康熙五十年（1711），康熙帝《御定佩文韵府》444卷，凡一万八千余页，也在扬州天宁寺刊刻告竣。档案〈奏为御颁《佩文韵府》在扬州刊刻工竣装箱进呈等请旨事〉（康熙五十二年九月初十日），记载此事。其文曰：

\begin{BookQuote}
臣李煦跪奏：窃臣煦与曹寅、孙文成，奉旨在扬州刊刻御颁《佩文韵府》一书，今已工竣。谨将连四纸刷钉十部，将乐纸刷钉十部，共装二十箱，恭进呈样，再连四纸应刷钉若干部、将乐纸应刷钉若干部，理合奏请，伏乞批示遵行，解送进京。臣煦临奏可胜悚惕之至。康熙五十二年九月初十日。(朱批）此书刻得好的极处。南方不必钉本，只刷印一千部，其中将乐纸二百
\end{BookQuote}

\vfill
\begin{BookNotes}
\BookNote{6.}{〈进全唐诗表〉，《全唐诗》，上海古籍出版社，上海，1986年。}
\BookNote{7.}{〈江宁织造曹寅奏谢列名全唐诗刊刻衔名折〉（康熙五十年三月初十日），中国第一历史档案馆编《康熙朝汉文殊批奏折汇编》，第3册，第375页，档案出版社，北京，1985年。}
\BookNote{8.}{彭定求，〈翰林院修撰东山汪君墓志铭〉，《四库全书存目丛书》，集部，第246册，第742页，齐鲁书社，济南，1997年。}
\BookNote{9.}{〈江宁织造曹寅奏报全唐诗集月内可以刻完折〉（康熙四十五年七月初一日），中国第一历史档案馆编《康熙朝汉文朱批奏折汇编》，第1册，第389页，档案出版社，北京，1985年。}
\BookNote{10.}{章宏伟，《十六—十九世纪中国出版研究》，第476页，上海人民出版社，上海，2011年。}
\BookNote{11.}{《四库全书总目》，第136卷，中华书局，北京，1965年。}
\end{BookNotes}

% 原书第 544 页（PDF 第 557 页）
\BookPage{544}{557}
部，即足矣。12

此期在扬州刊刻的第三部要籍是《御批资治通鉴纲目》。北宋司马光编纂《资治通鉴》，南宋朱熹则对其剪枝挺干、撰纲列目，而成《资治通鉴纲目》，便于简明阅读，益于资政治国。康熙帝说：司马光「编辑《资治通鉴》，论断古今，尽得其当」。13朱熹因读《资治通鉴》以作纲目，康熙帝因读《资治通鉴纲目》以作「御批」，而成《御批资治通鉴纲目》。康熙帝说自己仔细阅读「通鉴纲目」三遍，并对其析疑正漏，14折中权衡，亲加评定，御批着论，成《御批资治通鉴纲目》，凡59卷、前编1卷、外纪1卷、举要3卷、续编27卷。15这部《御批资治通鉴纲目》，也在扬州天宁寺刊刻竣工。

在清宫档案中，有《奏为遵旨刷印御批资治通鉴纲目〉事》折片，对此做出记载：

\begin{BookQuote}
管理苏州织造、大理寺卿、兼巡视两淮盐课监察御史臣李煦谨奏：恭请万岁万安。前臣煦遵奉圣旨，刷印《御批资治通鉴纲目》不必用套，草钉送来，钦此钦遵。臣煦鸠工刷钉六百部，从水路解运进呈，谨奏以闻。康熙四十九年三月
\end{BookQuote}

十九日。

\begin{BookQuote}
（朱批）知道了。16
\end{BookQuote}

上列康熙朝扬州文化中的刊刻《全唐诗》、《佩文韵符》和《御批资治通鉴纲目》三例盛事，以及康熙帝《御制诗集》在扬州刊刻，17说明扬州的雕印水准、精美装订和文化氛围，在清代文化史上占有重要地位，并为中华文化发展做出贡献。

在《清宫扬州御档选编》的康熙朝三十二件档案中，直接与文化有关的档案八

\vfill
\begin{BookNotes}
\BookNote{12.}{中国第一历史档案馆、扬州市档案局（馆）编:《清宫扬州御档选编》，第1册，第43页，广陵书社，扬州，2009年。}
\BookNote{13.}{《清圣祖仁皇帝庭训格言》，圣训第九十四条，清雍正八年，内府刻本，北京。}
\BookNote{14.}{《御批历代通鉴辑览·序》，卷首，清武英殿刻本，北京。}
\BookNote{15.}{《四库全书总目》，第755页，中华书局，北京，1965年。}
\BookNote{16.}{中国第一历史档案馆、扬州市档案局（馆）编：《清宫扬州御档选编》，第1册，第34页，广陵书社，扬州，2009年。}
\BookNote{17.}{中国第一历史档案馆、扬州市档案局（馆）编:《清宫扬州御档》，第1册，第227页，广陵书社，扬州，2009年。}
\end{BookNotes}

% 原书第 545 页（PDF 第 558 页）
\BookPage{545}{558}
\noindent 件，占本书收录档案的四分之一，18可见扬州文化的一斑。

还有，康熙朝纂修、雍正朝成书的《古今图书集成》，乾隆帝谕旨将其一部「发往扬州天宁寺行宫」19贮藏。

此外，嘉庆朝的《全唐文》一千卷，也在扬州进行慎校抽补、雕刻梨板、刊印装订、完竣呈送。「选编」中收录嘉庆朝有关《全唐文》的奏折四件：

其一、〈奏为承刊《钦定全唐文》在事诸臣及附近江浙官绅冀得捧函快诵请赏准各臣分印事〉（嘉庆二十一年十月二十七日）。先是，嘉庆十九年（1814）六月十九日，将内府《全唐文》写本一千卷、总目四卷，计一百函套，交两淮盐政在扬州刊校印装。20

其二、〈奏为遵旨校刊《钦定全唐文》完竣及装函进呈事〉（嘉庆二十一年十月二十七日）。奏折中称，锦套陈设本二十四部，每部装成一百套；石青杭细套赏赉本一百部，并日后将板片存贮运库。21

其三、〈奏为校刊《钦定全唐文》告成奉旨赏还顶戴谢恩事〉（嘉庆二十一年十二月十一日）。奏折请将该书赏给江浙等处官绅等，自备纸墨工价，分印赏赐。22

\vfill
\begin{BookNotes}
\BookNote{18.}{如：〈奏为御批高旻寺碑文事〉（康熙四十三年十二月初二日）、〈奏为僧人纪荫出任高旻寺主持事〉（康熙四十三年十二月初十日）、〈康熙帝御书「贤守清风匾额悬挂平山堂等事〉（康熙四十四年四月至闰四月）、〈奏为《全唐诗集》定期于天宁寺开局刊刻等事〉（康熙四十四年五月初一日）、〈奏为曹宜奉佛到扬州日期及前往普陀寺安置等事〉（康熙四十七年三月二十九日）、〈奏为遵旨刷印《御批资治通鉴纲目》事〉（康熙四十九年三月十九日）、〈奏为御颁《佩文韵府》在扬州刊刻工竣装箱进呈等请旨事〉（康熙五十二年九月初十日）、〈奏为遵旨斋戒择日在扬州天宁寺内延僧讽经力保皇太后圣体康宁事〉（康熙五十六年十一月初七日）。}
\BookNote{19.}{中国第一历史档案馆、扬州市档案局（馆）编：《清宫扬州御档选编》，第3册，第186页，广陵书社，扬州，2009年。}
\BookNote{20.}{〈奏为承刊《钦定全唐文》在事诸臣及附近江浙官绅冀得捧函快诵请赏准各臣分印事〉（嘉庆二十一年十月二十七日），中国第一历史档案馆、扬州市档案局（馆）编《清宫扬州御档选编》，第4册，第327页，广陵书社，扬州，2009年。}
\BookNote{21.}{〈奏为遵旨校刊《钦定全唐文》完竣及装函进呈事〉（嘉庆二十一年十月二十七日），中国第一历史档案馆、扬州市档案局（馆）编《清宫扬州御档选编》，第4册，第328页，广陵书社，扬州，2009年。}
\BookNote{22.}{〈奏为校刊《钦定全唐文》告成奉旨赏还顶戴谢恩事〉（嘉庆二十一年十二月十一日），中国第一历史档案馆、扬州市档案局（馆）编《清宫扬州御档选编》，第4册，第331页，广陵书社，扬州，2009年。}
\end{BookNotes}

% 原书第 546 页（PDF 第 559 页）
\BookPage{546}{559}
其四、〈奏为遵旨办理改装《全唐文》书函并分赴江浙文汇等阁安放事〉（嘉庆二十四年正月二十五日）。所云「改装」一事，系呈送的《全唐文》「每部一百函，每本皆有衬纸」，其函数较多，命将内府陈设之书各部，「撤去衬纸，每部改装五十函」；又在扬州文汇阁、镇江文宗阁、杭州文澜阁各贮存一部。23

扬州除刊刻诗书外，还有蒋衡书写经文盛事。

二

雍正朝蒋衡在扬州书写的「十三经」，后勒石竖碑，矗立于国子监，不仅是中华文化珍宝，而且是世界文化遗产。

雍正与扬州的文化关系，在《清宫扬州御档选编》的雍正朝六件档案中，直接与文化有关的档案有两件：其一是，〈遵旨覆奏拟以明纯为扬州高旻寺方丈并安排该员进京等事〉（雍正十二年四月初一日）；其二是，〈奏为扬州安定书院修理完竣并延师课训事〉（雍正十三年十月十八日）。以上两件关于「寺院」和「书院的折片，说明雍正帝对扬州文化的关注。然而，雍正朝扬州的一件文化盛事，就是蒋衡在扬州琼花观抄写「十三经」。

蒋衡（1672-1742），24江苏金坛人，初名衡，改名振生，字拙存，号湘帆，晚号拙叟人、拙老人。祖、父皆精书法，幼承家学，自小临摹，尤工行楷，苦练有成。科试落第后，便四出游学，寻师访友，切磋书艺，足迹半海内，尝「观碑关中，获晋、唐以来名迹，临摹三百余种」，刻成《拙存堂临古帖》28卷。

非历磨炼，难以大成。蒋衡科试不第，转意游学，研工书法。史书记载:「先生好远游，既不遇，遂东诣曲阜、谒孔林，至会稽，涉西江，历嵩少，导荆楚，登

\vfill
\begin{BookNotes}
\BookNote{23.}{〈奏为遵旨办理改装《全唐文》书函并分赴江浙文汇等阁安放事〉（嘉庆二十四年正月二十五日），中国第一历史档案馆、扬州市档案局（馆）编《清宫扬州御档选编》，第4册，第339页，广陵书社，扬州，2009年。}
\BookNote{24.}{蒋衡卒年，见有两说：其一，《清史列传·蒋衡》：「（乾隆）八年，卒，年七十二。」按：乾隆八年为1743年其二，王一宁着，〈清代书法大家蒋衡及其家族〉载述：「乾隆七年（I742）蒋衡病死，终年70岁。」（载《金坛文史选萃》上册，金坛市政协文史委员会编，2000年。)}
\end{BookNotes}

% 原书第 547 页（PDF 第 560 页）
\BookPage{547}{560}
\noindent 黄鹤矶，过大庾岭，升白鹤峰，访东坡故宅，抵琼海，观扶桑日出，登雁门山，历井陉，逾龙门，为终南岳之游，浴骊山温泉，登慈恩寺雁塔，纵观碑洞金石遗刻，所至以笔墨自随，赋诗作画，或歌哭相杂，至不能自止。」25在游学历程中，蒋衡长见识，展胸怀，摩碑刻，强意志，逐渐发愿要书写「十三经」。但他书写「十三经」是在长安碑林发愿的。

他在长安观摩碑林时，痛觉唐代《开成石经》出于众手，杂乱不齐，于是决心重写「十三经」一《周易》、《尚书》、《毛诗》、《周礼》、《礼记》、《仪礼》、《春秋左传》、《春秋公羊传》、《春秋谷梁传》、《论语》、《孟子》、《孝经》和《尔雅》。决心下定，矢志不移。雍正四年（1726）授英山教谕，力辞不赴。后又催促就职，仍以病老为由，上书求免。他在「邗江之琼花观」，26专心写经。「邗江」，古为邗沟，《左传》哀公九年:「秋，吴城邗，沟通江淮。」杜预注：「于邗江筑城穿沟。」27汉名邗沟为邗江。史载：秦置广陵县，隋改邗江县，元废。清人顾祖禹《读史方舆纪要》引《寰宇记》载:「在（扬州）西四里蜀岗上。」28今扬州市邗江区，以古邗江得名。29可见，「邗江之琼花观」其地在扬州。琼花观在扬州何处？在今扬州市文昌中路360号。30

琼花观，为扬州历史名胜之一。原为古后土祠（庙），建于西汉成帝元延二年（BC11）。北宋政和年间（1111-1118），徽宗赵佶赐额「蕃厘（xī）观」，遂改名为「蕃厘观」。后因观内有一株琼花，树茂花繁，洁白可爱，故又称「琼花观」。文人墨客，多有诗颂。宋臣韩琪诗曰：「维扬一株花，四海无同类。」刘敞诗则云:「东风万木竞纷华，天下无双独此花。」欧阳修任扬州知府时，在观内筑「无双亭」赏花，并诗曰:「琼花芍药世无伦，偶不题诗便怨人；曾向无双亭下醉，自知不负广陵春。」琼花美誉，驰名遐迩。31当年琼花观内，亭台楼榭，轩坊花

\vfill
\begin{BookNotes}
\BookNote{25.}{《国朝耆献类征·蒋衡传》初编，第433卷，光绪十六年（1890）刻本。}
\BookNote{26.}{《国朝耆献类征·蒋衡传》初编，第433卷，光绪十六年（1890）刻本。}
\BookNote{27.}{《左传》哀公九年，《十三经注疏附校勘记》，中华书局影印本，北京，1980年。}
\BookNote{28.}{顾祖禹，《读史方舆纪要》，卷23，上海书店出版社影印本，上海，1998年。}
\BookNote{29.}{《辞海》（第六版缩印本），「邗江」条，第696页，上海辞书出版社，上海，2010年。}
\BookNote{30.}{笔者在蒋鸿青先生陪同下，同爱新觉罗·启骧和李味辛夫妇考察琼花观遗址纪录。}
\BookNote{31.}{〈蕃厘观简介碑〉，碑在琼花观内，笔者抄录。}
\end{BookNotes}

% 原书第 548 页（PDF 第 561 页）
\BookPage{548}{561}
\noindent 石，几焚几建，遗韵犹存。今扬州以琼花为市花。这里在清雍正年间，成为蒋衡书写「十三经」的重要场所。

蒋衡在扬州琼花观，青灯相伴，中正灵静，32勤奋不辍，笃志写经。自雍正四年（1726）至乾隆二年（1737），历时12年，手抄「十三经」，楷书工整，六十余万字，终于大功告成。《清史稿·蒋衡传》记载:「键户十二年，写十三经。乾隆中，进上，高宗命刻石国学，授衡国子监学正，终不出。」33

大成垂名，常在身后。蒋衡所书「十三经」，身后五十年，乾隆帝前命将蒋衡所书写的「十三经」刻石，贞珉工竣，御制序文，立于太学，以垂万世。

蒋衡书写「十三经」的成功，得到三位贵人的鼎助：

第一位是扬州富商马日琯。马氏出资两千金，将其手书「十三经」装裱成册。史载:「扬州马日琯为出白金二千锾，装潢成三百册，五十函。」34

第二位是江南河道总督高斌。乾隆五年（1740），高斌将蒋衡手书的「十三经」上奏:「总河高斌奏称，江南镇江府金坛县贡生蒋振生，依石经式，手书『十三经』正文，计三百册，共五十函。谨先进《易经》二函，可否将全册五十函进呈？」得旨进呈后，经武英殿儒臣详加校阅，奏报的结论是:其《易经》两册，经审查，「字画尚属端楷」。尔后，得旨:「令高斌将全册五十函，送交武英殿，再加校定。」又经审定:「经文果无讹误，字画一律端好。」于是，「大学士等以该生年近七旬，志在尊经，请赏给国子监学正职衔。其手书『十三经』，请用枣木板镌刻刷印，以备颁发。疏入，报闻」35。

第三位是乾隆皇帝。蒋衡手写「十三经」进呈后，乾隆帝先要将其雕版印刷，但受阻未果。于是，乾隆五十六年（1791），命以蒋衡手书「十三经」为底本，刻石太学，定名《乾隆石经》。乾隆五十九年（I794)，石碑镌刻成，立于国子监。

\vfill
\begin{BookNotes}
\BookNote{32.}{蒋衡，〈书法论〉，见《国朝耆献类征》第433卷，第53页上，光绪十六年（1890）刻本。}
\BookNote{33.}{《清史稿·蒋衡传》，第503卷，列传二百九十第13888页，中华书局标点本，1977年。}
\BookNote{34.}{《清史列传》，第71卷，〈文苑传〉二，页5856，中华书局标点本，1987年。}
\BookNote{35.}{《清高宗纯皇帝实录》，卷99、乾隆四年八月庚寅十六日，中华书局影印本，北京，1985年。}
\end{BookNotes}

% 原书第 549 页（PDF 第 562 页）
\BookPage{549}{562}
\noindent 全部刻石碑189通，加「谕旨告成表文碑1通，共190通，现藏于北京孔庙和国子监博物馆内。

蒋衡手书、乾隆刻石的「十三经刻石」，即「乾隆石经」，其规模之宏大，楷法之工整，笔力之雄健，毅力之坚韧，学志之专一，价值之珍贵，国内仅有，世界也无，从而成为中国、也成为世界文化艺术宝库中的稀世珍品。

此外，今存扬州大明寺东墙外之「淮东第一观」五字，每字约一米见方，笔力遒劲，挥洒廓大，为蒋衡所手书。

蒋衡书写《十三经》时一度寓居扬州，死后葬于扬州大明寺外斜坡下。他死后，扬州文化在乾隆朝又有新的发展。

三

乾隆帝继位之后，扬州文化之地位、之贡献，远超康熙和雍正时期。在《清宫扬州御档选编》的乾隆朝一百零二件档案中，直接与文化有关的档案四十五件，约占其总数的百分之四十四。其中与书籍等相关的档案二十一件，占总数的四分之一强。而与《四库全书》有关的档案十件，约占同书相关档案的近百分之五十。乾隆朝扬州文化的一个特点，是同纂修《四库全书》密切关联。这主要表现在两个方面:一方面是，《四库全书》原本大量来自扬州；另一方面是，《四库全书》藏书七阁之一为扬州文汇阁。

扬州为《四库全书》的重要书源之一。先是，扬州有怀素真迹36、米芾手卷37和苏轼墨迹3等珍贵文物进献，在编修《四库全书》征集书籍时，乾隆帝自然想到扬

\vfill
\begin{BookNotes}
\BookNote{36.}{〈奏为遵旨觅得藏经纸等物并将家传怀素等等真迹恭呈御览事〉（乾隆三十五年五月二十六日），载中国第一历史档案馆、扬州市档案局（馆）编《清宫扬州御档选编》，第2册，第147页，广陵书社，扬州，2009年。}
\BookNote{37.}{〈奏为遵旨觅得米芾手卷等物于装修船上附带进京并请免开价值事〉（乾隆三十五年闰五月初二日），载中国第一历史档案馆、扬州市档案局（馆）编《清宫扬州御档选编》，第2册，第148页，广陵书社，扬州，2009年。}
\BookNote{38.}{〈奏为觅得宋苏轼墨迹手卷等恭呈御览事〉（乾隆三十五年九月初三日），载中国第一历史档案馆、扬州市档案局（馆）编《清宫扬州御档选编》，第2册，第149页，广陵书社，扬州，2009年。}
\end{BookNotes}

% 原书第 550 页（PDF 第 563 页）
\BookPage{550}{563}
\noindent 州

\begin{BookQuote}
《四库全书》之原书大量来自扬州。「选编」中，收录扬州等藏书御档七件，
\end{BookQuote}

如〈奏为采选书籍二百九十一种派专差赍送事〉（乾隆三十八年闰三月三十日）、〈扬州商人马裕家藏好书挑交四库全书总裁事〉（乾隆三十八年四月十九日）、〈上谕编四库全书征集民间旧书事〉（乾隆三十八年五月十七日）、〈奏为已得书籍专差赍送事〉（乾隆三十八年六月初八日）、〈奏为续得书籍专差赍送事〉（乾隆三十八年六月二十四日）、〈奏为恭进书籍事〉（乾隆三十八年七月十一日）和〈奏为续购各书专差赍送事〉（乾隆三十八年九月二十八日）等。仅以扬州商人马裕为例。

扬州商人马裕家里藏书丰富，乾隆帝有所耳闻:「上谕:前以办理《四库全书》，闻扬州商人马姓家内藏书颇富，曾传谕两淮盐政李质颖，令其就近妥协访问借抄。」39其书为借抄，用毕应发还。但有些书用完存库，并未归还。后乾隆帝谕旨:「所有底本，将来俱应存贮翰林院衙门。」40马裕家藏书，经过核查，共一千三百八十五种。

朝廷对马裕家藏书有个了解过程，马裕对朝廷「借书」也有个信任过程。以「选编」所载四件御档为例:乾隆三十八年四月十九日，两淮盐政李质颖奏报，马裕家先检出一百三十三种，后又检出六十二种，共计一百九十五种。但乾隆帝认为马裕「未免心存畏惧，又惮将善本远借」，故乾隆帝在上谕中强调将书「上紧抄进，仍将原本发还」，并令李质颖「善为寻觅」。41尔后，马裕进呈书目，选取二百一十种，督臣高晋续选六十二种，又拣出三百七十种。再后将马裕家一千三百八十五种书籍，选取七百七十六种呈送。到同年五月十七日，内阁奉上谕，江浙督抚及两淮盐政共购求书籍四五千种，而马裕一家拣选之书约占江浙督抚

\vfill
\begin{BookNotes}
\BookNote{39.}{中国第一历史档案馆、扬州市档案局（馆）编：《清宫扬州御档选编》，第3册，第164页，广陵书社，扬州，2009年。}
\BookNote{40.}{《清高宗纯皇帝实录》卷1282、乾隆五十二年六月戊申十二日，中华书局影印本，北京，1985年。}
\BookNote{41.}{中国第一历史档案馆、扬州市档案局（馆）编:《清宫扬州御档选编》，第3册，第165页，广陵书社，扬州，2009年。}
\end{BookNotes}

% 原书第 551 页（PDF 第 564 页）
\BookPage{551}{564}
\noindent 及两淮盐政购求书籍总数的近百分之二十，数量之大，殊为可观。42

\begin{BookQuote}
《四库全书》不仅书源之一在扬州，而且书藏之一也在扬州。
\end{BookQuote}

扬州文汇阁为《四库全书》的七阁藏书之一。《四库全书》七分写本陆续告竣之后，分别贮藏于北四阁一紫禁城文渊阁、圆明园文源阁、避暑山庄文津阁和盛京沈阳文溯阁，南三阁一扬州文汇阁、镇江文宗阁和杭州文澜阁。43杭州为浙江省会，也是杭州将军驻地；镇江在江南；而扬州既不在江南，也并非省会，更不是将军驻地一江宁（南京）为江宁将军驻地、六朝古都，苏州为江苏巡抚驻地、人文荟萃，但《四库全书》并未分藏在江宁和苏州，却分藏在江北的扬州，当别有一番深意。

\begin{BookQuote}
《四库全书》南三阁之一为扬州文汇阁。「选编」收录档案三件，即〈奏为
\end{BookQuote}

扬州天宁寺等地藏书楼盖造完竣请旨颁赐御书匾额事〉（乾隆四十四年十月二十六日）、〈上谕《四库全书》告竣分贮扬州大观堂之文汇阁等处事〉（乾隆四十七年七月初八日）和〈奏为遵旨交办文宗文汇二阁书籍事〉（乾隆五十五年十一月初九日）。

到乾隆四十七（1782）年七月初八日，《四库全书》的头分已经告竣，其第二、第三、第四分，限于六年之内，按期抄写完毕，分贮于文渊阁、文溯阁、文源阁和文津阁。乾隆五十五年（1790），「上谕:《四库全书》现在头分已经告竣，其二三四分，限于六年内，按期蔵事，并特建文渊、文溯、文源、文津等阁，以供藏庋。因思江浙为人文渊薮，允宜广布流传，以光文治，现特发内帑银，雇觅书手，再行缮写全书三分，分贮杨（扬）州大观堂之文汇阁，镇江金山寺之文宗阁，杭州圣因寺内拟改建文澜阁，以昭美备。」44其扬州大观堂已经贮藏《古今图书集成》，命在其空格里收储《四库全书》，如不敷用，再行添补。着闽浙总督兼浙江

\vfill
\begin{BookNotes}
\BookNote{42.}{中国第一历史档案馆、扬州市档案局（馆）编：《清宫扬州御档选编》，第3册，第167页，广陵书社，扬州，2009年。}
\BookNote{43.}{《四库全书》底本藏于翰林院。见《清高宗实录》卷1355、乾隆五十五年五月癸卯二十三日。但有学者以未见其实物传存而置疑之。}
\BookNote{44.}{中国第一历史档案馆、扬州市档案局（馆）编:《清宫扬州御档选编》，第3册，第207页，广陵书社，扬州，2009年。}
\end{BookNotes}

% 原书第 552 页（PDF 第 565 页）
\BookPage{552}{565}
\noindent 巡抚陈辉祖等督办，其费用由两浙商人捐办。45

在办理过程中，发现存在问题。乾隆五十二年（1787）六月，乾隆帝谕道：「朕前披阅文津阁所贮《四库全书》，看到其中讹谬甚多。令大学士将文渊、文源二阁所贮书籍一体校阅，其讹舛处不一而足。如阎若璩《尚书古文疏证》一书，有引李清、钱谦益诸说未经删削，并有连篇空白页，实属草率已极。着将承办之总校、分校等，交部议处。凡有违碍，即行修改，再行赔写抽换，务期完善等语。并命武英殿提调、总校陆费墀对文澜、文汇、文宗三阁三分「所有面页、装钉、木匣、刻字等项，俱着陆费墀自出己赀，仿照文渊等三阁式样罚赔，妥协办理。」46《四库全书》南三阁的三分全书，到乾隆五十五年（1790）才算告竣。

陆费墀时已身故，事情却有了变化。乾隆五十五年（1790）十月，「谕军机大臣等，前因江、浙两省，为人文之薮，特将《四库全书》添办三分，发交扬州、金山及杭州文宗、文汇、文澜三阁藏贮，所有装潢庋架等事，俱交两淮、浙江盐政办理。嗣因陆费墀总理《四库全书》，草率错误，获咎甚重，即罚令出赀承办。陆费墀本系寒士，家无担石，向在于敏中处，藉馆为业，谅不过千金产业耳！今所办三阁书匣等项，及缴出罚银一万两，计其家赀已不下二三万，若非从前在四库馆提调任内，苞苴馈送，何以有此多赀？现在陆费墀业已身故，所有插架、装匣等事，若令伊子接办，恐未能谙习。且身后所遗家业，想已无多，亦难措办。此时三分书俱已校对完竣，自应全行发往三处藏庋，未便稽延。着传谕海甯、全德，即仿照前次发去装潢、书匣等式样制造，专派妥商办理。并着海甯查明陆费墀原籍现有田房产业，加恩酌留一千两之数，为伊家属养赡，如尚有余赀，即作为添补三阁办书之用。海甯、全德务须认真督率该商等经理，妥速蔵工，毋任迟延草率。」47

乾隆五十五年（1790）十一月初九日，《四库全书》扬州、镇江两分已经抄写

\vfill
\begin{BookNotes}
\BookNote{45.}{中国第一历史档案馆、扬州市档案局（馆）编：《清宫扬州御档选编》，第3册，第208页，广陵书社，扬州，2009年。}
\BookNote{46.}{《清高宗纯皇帝实录》卷1282、乾隆五十二年六月戊申十二日，中华书局影印本，北京，1985年。}
\BookNote{47.}{《清高宗纯皇帝实录》卷1365、乾隆五十五年十月戊辰二十一日，中华书局影印本，北京，1985年。}
\end{BookNotes}

% 原书第 553 页（PDF 第 566 页）
\BookPage{553}{566}
\noindent 告竣，并已校对，其装订、装潢、书匣、书架，按照内府所发式样，由地方办理。48此事，《清高宗实录》记载:「谕:《四库全书》荟萃古今载籍，至为美备。不特内府珍藏，借资乙览，亦欲以流传广播，沾漑艺林。前因卷页浩繁，中多舛错，特令总纂等复加详细讐校，俾无鲁鱼亥豕之讹。兹已厘订蔵工，悉臻完善。所有江、浙两省文宗、文汇、文澜三阁，应贮全书，现在陆续颁发藏庋。该处为人文渊薮，嗜古好学之士，自必群思博览，藉广见闻。从前曾经降旨，准其赴阁检视钞录，俾资搜讨。但地方有司，恐士子等翻阅污损，过为珍秘，以阻其争先快睹之忱。则所颁三分全书，亦仅束之高阁，转非朕搜辑全书津逮誉髦之意。即武英殿聚珍版诸书，排印无多，恐士子等亦未能全行购觅。着该督抚等，谆饬所属，俟贮阁全书排架齐集后，谕令该省士子有愿读中秘书者，许其呈明，到阁钞阅，但不得任其私自携归，以致稍有遗失。…互为钞录传之日久，使石渠天禄之藏，无不家弦户诵，益昭右文稽古嘉惠士林盛事，不亦善乎！」49

事过四年，又起波澜。《清高宗实录》记载，乾隆五十九年（1794)七月，谕曰:「朕披阅《通鉴辑览》内，唐开元五年九月，令史官随宰相入侍，群臣对仗奏事条下，引贞观旧制，诸司皆正邪奏事。又注称:唐大明宫含元殿为正邪，亦谓之南邪等语。心疑笔误，因查诸旧史，乃知俱将牙字误作邪字。更检阅字书，牙与衙字，本属通用。至邪字从无与牙字相通之义，甚为舛错。所有原办《通鉴辑览》之总裁、纂修、校对等官，现经军机大臣查明，原书首卷开列衔名内，现存各员，如阿桂、孙士毅、纪昀、彭元瑞、毕沅、吴省钦等，本应交部议处，姑念成书已久，事隔多年，阿桂等及其余纂修校对各官，俱着一体加恩，免其交部。至武英殿刊本，及《四库全书》缮本，俱着查明改正外，所有颁行各直省刻本，并盛京、江、浙省文溯、文宗、文汇、文澜四阁存贮缮本，亦着各该督抚府尹等，一律改正。」50

贮藏《四库全书》的文汇阁，在扬州大观堂旁侧。据《扬州画舫录》云「御书楼在御花园中，园之正殿名大观堂，楼在大观堂之旁，恭贮《颁定图书集成》全

\vfill
\begin{BookNotes}
\BookNote{48.}{〈奏为遵旨交办文宗文汇二阁书籍事〉（乾隆五十五年十一月初九日），载中国第一历史档案馆、扬州市档案局（馆）编《清宫扬州御档选编》，第3册，第327页，广陵书社，扬州，2009年。}
\BookNote{49.}{《清高宗纯皇帝实录》卷1355，乾隆五十五年五月癸卯二十三日，中华书局影印本，北京，1985年。}
\BookNote{50.}{《清高宗纯皇帝实录》卷1456，乾隆五十九年七月乙未初十日，中华书局影印本，北京，1985年。}
\end{BookNotes}

% 原书第 554 页（PDF 第 567 页）
\BookPage{554}{567}
\noindent 部，赐名文汇阁，并『东壁流辉』匾。…文汇阁凡三层，忙瘤楹柱之间，俱绘以书卷。最下一层，中供《图书集成》，书面用黄色绢；两畔橱皆经部，书面用绿色绢；中一层，尽史部，书面用红色绢；上一层，左子右集，子部书面用玉色绢，集用藕合色绢。其书帙多者，用楠木作函贮之；其一本二本者，用楠木版一片夹之，束之以带，带上有环，结之使牢。」51

文汇阁本《四库全书》的数量，《续金山志》载文宗阁本数量，虽不尽同，却资参考。其数量为:经部九百四十七匣，五千四百零二本；史部一千六百二十五匣，九千四百六十三本；子部一千五百八十三匣，九千零八十四本；集部二千零四十二匣，一万二千三百九十八本。总目录计二十二匣，一百二十七本。52

但是，时间仅过五十多年，扬州文汇阁所藏《四库全书》便遭厄难。

扬州文汇阁贮藏的《四库全书》，毁于战火。在一份〈扬州阖属士民公告启〉中说:「敬启者:自二月二十三日，贼匪陷城，所有绅士百姓，逃出者居半，被陷者居半一其逃出者虽弃其家，尙存其命；而坐陷者，或悬梁、或投井、或自焚、或被掳，全家覆没者，不知凡几，惨何可言！贼匪始而搜掳，既而焚烧。两月以来，城内外殆将烧尽。昔日之花团锦簇，今日之断井颓垣，真有目不忍睹，耳不忍闻者！」53在这场战火中，「扬城已为灰烬」。54扬州城的建筑与文物，遭到焚毁，甚为可叹。55文汇阁及其贮藏的《四库全书》，焚烧殆尽，殊为可惜！

当时繁华城市有「两京两州」即北京、南京和扬州、苏州之誉。在两京、两

\vfill
\begin{BookNotes}
\BookNote{51.}{李斗撰，《扬州画舫录》，第4卷，第103-104页，中华书局，北京，1960年。}
\BookNote{52.}{参见吴哲夫《四库全书纂修之研究》，第148页，台湾故宫博物院印，台北，1990年。}
\BookNote{53.}{见〈呈贼匪陷毁扬州城早将误国害民之帅撤回阖属士民公告启〉（咸丰三年五月十二日），载中国第一历史档案馆、扬州市档案局（馆）编《清宫扬州御档选编》第3册第414页，广陵书社，2010年；又见〈奏报贼匪陷毁扬州城上达天聪早将误国害民之帅撤回阖属士民公告单〉（咸丰三年五月十二日），载中国第一历史档案馆、扬州市档案局（馆）编《清宫扬州御档》，第17册，第11703页，广陵书社，扬州，2009年。}
\BookNote{54.}{奏报扬州城战况记事单〉（咸丰三年五月十二日），载中国第一历史档案馆、扬州市档案局(馆)编《清宫扬州御档》，第17册，第11703页，广陵书社，扬州，2010年。}
\BookNote{55.}{《清宫扬州御档选编》与《清宫扬州御档》两书所录之两份档案，虽内容雷同，但标题不同。如上引《清宫扬州御档》作〈奏报扬州城战况记事单〉（咸丰三年五月十二日）；《清宫扬州御档选编》却作〈呈扬州城战况记事单（咸丰三年五月十二日）等。}
\end{BookNotes}

% 原书第 555 页（PDF 第 568 页）
\BookPage{555}{568}
\noindent 州中，选在扬州设诗局刻印《全唐诗》，同盛清扬州文化有关。盛清时期的扬州文化，之所以发达，主要是由于一其一，舆地因素：濒临长江、淮河、黄河、运河，东为海，系水陆交通枢纽，扼南北漕运咽喉；其二，历史因素：近二千五百年文化积累，尤其是隋唐、明代和清朝（前期）的三次繁盛；其三，政治因素：康熙皇帝和乾隆皇帝都曾六次南巡，驻跸扬州；其四，经济因素：两淮盐商，财货积厚，捐纳重金，赞助书籍出版；其五，人文因素：人文荟萃，文化昌盛；其六，工艺因素:雕版印刷，技艺高超。所印之书，康熙帝赞誉为「此书刻得好的极处」。

总之，盛清时期扬州文化，已经成为历史。然而，综观盛清扬州文化，仅从《全唐诗》、「十三经刻石」和《四库全书》三个文化实例，可以看出盛清社会与扬州文化的互动关系一扬州文化既拥簇盛清文化的繁荣，盛清文化也催促扬州文化的繁华；扬州文化为中华文化历史的发展和繁荣做出了自己的贡献。


% 原书第 556 页（PDF 第 569 页）
\BookPage{556}{569}
